\paragraph{A positive mass: the class of planes fixed by the regulator (algebraic at every \(N\), proved at \(N=2\)).}
Keep the setting and notation of the two preceding paragraphs, with Definition~(A) and the mass term, and let \(\gamma>0\). Write \(\epsilon+\mathrm i\gamma'=|\epsilon+\mathrm i\gamma'|\mathrm e^{\mathrm i\varphi}\) as in (i) of the preceding paragraph, now with \(0<\varphi<\pi/2\), and put \(h_\varphi=\mathrm e^{\mathrm i\varphi}h\). Consider paths \(\Lambda(\tau)\), \(0\le\tau\le1\), piecewise continuously differentiable, with \(\Lambda(0)\) real, \(\operatorname{Re}\mathrm iS_{\mathrm b}\le0\) on every contour of the path and \(\operatorname{Re}M_\gamma(\tau)\) positive definite, the conditions used for the Stokes arguments above.

(i) \emph{Classes of planes.} With \(\ell=a+\mathrm ib\) as in the preceding paragraph, the identity \((\epsilon+\mathrm i\gamma')h=M_\gamma-2\epsilon\ell\ell^{\mathrm T}\) shows that \(\operatorname{Re}h_\varphi\) is positive definite on the real subspace \(a^{\mathrm T}x=b^{\mathrm T}x=0\), of dimension at least eight, whenever \(\operatorname{Re}M_\gamma\) is positive definite. Let \(E\) be a real 2-plane on which \(\operatorname{Re}h_\varphi\) is positive definite, let \(r\) and \(s\) be the matrices of \(\operatorname{Re}h_\varphi\) and \(\operatorname{Im}h_\varphi\) on \(E\) in a real basis, and let \(\sigma_1,\sigma_2\) be the eigenvalues of \(r^{-1}s\), which are real and do not depend on the basis. Then \(\Theta(E)=\arctan\sigma_1+\arctan\sigma_2\in(-\pi,\pi)\) is a continuous branch of \(\arg\det(h_\varphi|_E)\). At \(N=2\) the condition \(\operatorname{Re}\mathrm iS_{\mathrm b}\le0\) is \(\operatorname{Im}[h(a',a')h(b',b')-h(a',b')^2]\ge0\) for all real \(a',b'\), and for \(N\ge3\) it implies this, by restriction to an \(\mathfrak{su}(2)\) block. Since \(\det(h|_E)=\mathrm e^{-2\mathrm i\varphi}\det(h_\varphi|_E)\) in the same basis, it gives \(\sin(\Theta(E)-2\varphi)\ge0\). So \(\Theta(E)\) lies either in \([2\varphi,\pi)\), the upper class, or in \((-\pi,2\varphi-\pi]\), the lower class, and these are separated by an interval of length \(\pi\).

(ii) \emph{The class is conserved, and \(-\mathfrak C_0\) is not reached.} For a real symmetric form with at least two positive eigenvalues the planes on which it is positive definite form a connected set: setting the components of each vector in the eigenspaces of non-positive eigenvalues to zero, by a factor decreasing from \(1\) to \(0\), keeps the form positive definite on the plane, and the planes in the positive eigenspace form a connected Grassmannian. Along a path as above \(\operatorname{Re}h_\varphi\) has at least eight positive eigenvalues by (i), so all planes on which it is positive definite have one class at each \(\tau\), and by continuity of \(\Theta\) this class does not depend on \(\tau\). At \(\tau=0\), \(h\) is real with the inertia of \(\eta\), so \(\Theta(E)=2\varphi\) and the class is upper. Now let \(h\in-\mathfrak C_0\) with \(\operatorname{Re}M_\gamma\) positive definite, and write \(h=-T^{\mathrm T}\operatorname{diag}(d_\mu)T\) with the arguments \(\beta_\mu\) of \(\mathfrak C_0\), at least eight of which lie in \((0,\pi/2)\) by (ii) of the paragraph on linear contours that mix the indices. Since \(\operatorname{Re}h_\varphi=-T^{\mathrm T}\operatorname{diag}(\operatorname{Re}\mathrm e^{\mathrm i\varphi}d_\mu)T\), (i) and the law of inertia give \(-\operatorname{Re}\mathrm e^{\mathrm i\varphi}d_\mu>0\) for at least eight \(\mu\). So there are \(\mu\neq\nu\) with both properties; for them \(\beta_\mu,\beta_\nu>\pi/2-\varphi\), the plane spanned by \(T^{-1}e_\mu\) and \(T^{-1}e_\nu\) is positive, and \(\Theta=\beta_\mu+\beta_\nu+2\varphi-2\pi<2\varphi-\pi\), the lower class. Hence, for every \(N\), no path as above joins the Hermitian contour to a contour with \(h\in-\mathfrak C_0\). In particular it does not reach the contour \(A_0=\mathrm e^{\mathrm i\pi/8}X_0\), \(A_i=\mathrm e^{-3\pi\mathrm i/8}X_i\) of (i) of the preceding paragraph. For \(\gamma\le0\), where \(-\pi/2\le\varphi\le0\), the interval \((-\pi,2\varphi-\pi]\) is empty, and the inertia argument of (i) of the preceding paragraph applies instead. The same proof holds for the regulator \(\epsilon\sum_{\mu,\nu}P_{\mu\nu}\operatorname{Tr}A_\mu A_\nu\), with \(P\) positive definite and \(P^{-1}\eta\) having a nine-fold eigenvalue \(1/c\), with \(\epsilon c\) in place of \(\epsilon\) in \(\varphi\), since \(P-c\eta\) then has rank one.

(iii) \emph{The limit \(\epsilon\to0\).} Let \(\epsilon_j\to0\) and \(\gamma_j>0\) with \(\gamma_j/\epsilon_j\to\infty\), which includes fixed \(\gamma>0\), and let \(\Lambda_j\) be the end points of paths as above for \((\epsilon_j,\gamma_j)\), converging to an invertible \(\Lambda_\infty\) with \(h=h_\infty\). Then \(\varphi_j\to\pi/2\) and \(\operatorname{Re}M_{\gamma_j}/|\epsilon_j+\mathrm i\gamma_j'|\to-\operatorname{Im}h_\infty\), with \(\epsilon_jc\) in place of \(\epsilon_j\) for the regulator of (ii), so \(-\operatorname{Im}h_\infty\) is positive semidefinite. If it had two positive eigenvalues, a plane on which it is positive definite would, for large \(j\), be positive for \(\operatorname{Re}h_{\varphi_j}\) and of the upper class, with \(\Theta\ge2\varphi_j\to\pi\), whereas \(\Theta\) on it converges to a value in \((-\pi,\pi)\). So \(\operatorname{Im}h_\infty=-\kappa\,vv^{\mathrm T}\) with \(\kappa\ge0\) and \(v\) real. On the configurations with \(\sum_\alpha v_\alpha X_\alpha=0\), \(S_{\mathrm b}\) and the mass term are then real at every \(N\), and at \(\epsilon=0\) the integrand has modulus \(|\operatorname{Pf}\mathcal M|\) there. At \(N=2\) this integrand of \(I[1]\) on the contour \(\Lambda_\infty\), with any real \(\gamma\) in the mass term, is not absolutely integrable. The restriction of \(h_\infty\) to \(v^\perp\) is real of rank at least eight, so \(|\det\mathcal G|\ge c'R^6\) on a cone of colour columns in \(v^\perp\) of norm \(R\). Shifting the columns along \(v\) by at most \(R^{-3}\) changes the real part of the exponent by a bounded amount and \(\det\mathcal G\) by \(O(R^2)\), so the modulus integrates to at least a constant times \(\int^\infty R^{26}R^{24}R^{-9}\,\mathrm dR=\infty\). So for fixed \(\gamma>0\), and for \(\gamma/\epsilon\to\infty\) with \(\gamma\) converging to a finite limit, no linear contour reached from the Hermitian contour in this way, fixed or convergent as \(\epsilon\to0\), carries the dominated convergence used above for \(\gamma\le0\), since a dominating function would also dominate the integrand at \(\epsilon=0\) on \(\Lambda_\infty\) with that limit of \(\gamma\); this is in line with the closing of the admissible \(u\) on \(\mathcal R_u\) like \(\epsilon/\gamma\).

So, for linear contours, Definition~(A) or a regulator as in (ii), fixed \(\gamma>0\) and \(\epsilon\to0\) taken at fixed \(N\), the contour of \(-\mathfrak C_0\) found in the preceding paragraph is not joined to the Hermitian contour, and at \(N=2\) no contour joined to it, fixed or convergent as \(\epsilon\to0\), has a limit on which the integrand at \(\epsilon=0\) is absolutely integrable. The mass is a supplied deformation of the action here. Whether the limit at fixed \(\gamma>0\) exists, whether it differs from the Euclidean model, and what follows when \(\gamma\to0\) after it, are not decided; for \(\gamma\to0\) first, or with \(\gamma/\epsilon\) bounded, the paragraph on the regulated Lorentzian integral gives the Euclidean values at phases. This is no statement about non-linear contours, about linear maps that also mix the colour components, about paths on which \(\operatorname{Re}M_\gamma\) is only positive semidefinite, about other positive definite regulators, about families that do not converge, about \(N\to\infty\) taken first, or about Definitions~(B) and~(C). The two classes are fixed by the \(\eta\) in \(S_{\mathrm b}\) and by the regulator, so they select no signature, and the positive mass does not discharge Assumption~\ref{ass:lorentzian}.
